% !TeX root = ../../Book3.tex
\chapter{素数定理}
\label{b3:ch:27}
\BookChapterNumber{b3:ch:27}

\begin{chapterabstract}
本章把 zeta 函数的边界解析性质转化为素数的计数渐近。先建立不依赖素数定理的 Chebyshev 双边估计，说明普通计数、对数权和与素数幂权和的关系；再用非负三角多项式排除实部一上的零点。主线完整证明解析边界版本的 Wiener--Ikehara 定理，包括累积函数的测度表示、正核平均及单调夹逼。选读部分另用 Newman 的围道方法完成同一结论。两条证明均给出定性主项，并明确指出进一步估计误差还需要什么资料。
\end{chapterabstract}

\begin{learninggoals}
能够解释为什么解析证明先处理 \(\psi\) 或 \(\vartheta\)，再恢复 \(\pi\)；能够在无零点论证中追踪非负权重、三角多项式与局部阶数的作用。能够恢复 Tauber 变换中的变量与归一化，检查整条边界延拓、端点质量、核的频率支集和极限次序；选读后能够区分 Newman 证明中允许用于 \(g\) 与截断整函数 \(g_T\) 的两种围道。能够把素数定理用于固定比例区间和加权计数，同时辨认短区间、次项及有效误差的不同要求。
\end{learninggoals}

素数定理断言
\[
 \pi(x)\sim\frac{x}{\log x}\qquad(x\to\infty),
\]
其中 \(\pi(x)\) 是不超过 \(x\) 的素数个数。右边没有给出每个整数是否为素数的判据，而是描述一个越来越长的初始区间内的累计数量。即使已经从初等方法知道 \(\pi(x)\) 被两个正常数倍的 \(x/\log x\) 夹住，仍需证明这两个常数能够收紧到同一个数一。

前两章建立了算术与复分析的联系。\chapref{b3:ch:25}将算术函数写成 Dirichlet 级数，并由 Euler 乘积得到
\[
 -\frac{\zeta'(s)}{\zeta(s)}
 =\sum_{n\ge1}\Lambda(n)n^{-s},\qquad \operatorname{Re}s>1.
\]
\chapref{b3:ch:26}则证明 \(\zeta\) 在全平面亚纯，并在一处有留数为一的单极点。但这两项资料还不是素数定理：一处的极点只描述一种平均发散，尚不能排除其他边界奇性，也不能独自保证原计数函数有渐近。

本章将这两个缺口分开处理。首先，非负系数与恒等式
\(3+4\cos u+\cos2u=2(1+\cos u)^2\)
排除 \(\operatorname{Re}s=1\) 上的零点，使去掉主部的对数导数能够越过整条边界。其次，Tauber 方法把这项解析信息转成累积量的渐近；其中单调性用来防止局部偏离藏在平均之下。无零点性和单调性解决的是不同问题，不能互相替代。

基本阅读路线为\secref{b3:sec:2701}至\secref{b3:sec:2705}。\secref{b3:sec:2704}先恢复一般右连续累积函数的 Stieltjes 测度和 Laplace 变换，再用正核完整证明 Wiener--Ikehara 的解析边界版本。\secref{b3:sec:2706}为选读，它可从前三节直接进入，通过一个较短的围道引理重新证明素数定理，不以第四、五节的结论为前提。本章两条路线都不依赖上一章选读的 \(\xi\) 乘积或无穷零点显式公式。

Newman 的《Simple Analytic Proof of the Prime Number Theorem》\cite{Newman1980Simple}给出末节方法的历史出处；本章实际展开的围道论证依 Zagier 的《Newman's Short Proof of the Prime Number Theorem》\cite[pp. 706--708]{Zagier1997Newman}核对。阅读这类短证明时，尤其要补回有限高度与全边界的区别，以及先固定围道、再扩大围道的极限次序。这里将这些容易被篇幅压缩的步骤写在相应计算旁边。

% !TeX root = ../../Book3.tex
\section{素数计数函数}
\label{b3:sec:2701}

素数定理首先是一个计数断言。解析方法直接控制的是对数权和及素数幂的和，因此需要先说明三种计数函数为何具有等价的首项渐近，并建立一个独立于定理的粗上界。

\subsection{\texorpdfstring{\(\pi(x)\)}{\AnalysisPdfPi(x)}}
\label{b3:subsec:270101}

记 \(\pi(x)=\#\{p\le x:p\text{ 为素数}\}\)，\(x\ge1\)。本章始终包含端点处的素数，使函数右连续。若改用严格小于，差至多一，不改变最终渐近，但会影响精确的 Stieltjes 积分端点约定。

\subsection{\texorpdfstring{\(\vartheta(x)\)}{\AnalysisPdfVarTheta(x)}}
\label{b3:subsec:270102}

定义 \(\vartheta(x)=\sum_{p\le x}\log p\)。乘法形式为 \(e^{\vartheta(x)}=\prod_{p\le x}p\)。
\begin{lemma}[Chebyshev 上界]\label{b3:lem:27-chebyshev}
有 \(\vartheta(x)=O(x)\)。
\end{lemma}
\begin{proof}
每个 \(n<p\le2n\) 的素数整除 \(\binom{2n}{n}\)，所以
\(\vartheta(2n)-\vartheta(n)\le\log\binom{2n}{n}\le2n\log2\)。对任意实数 \(x\)，取邻近整数处理端点，得 \(\vartheta(x)-\vartheta(x/2)\le Cx\)：长度至多二的剩余区间至多含两个整数，各贡献不超过 \(\log(x+2)\)。沿 \(x,x/2,x/4,\ldots\) 累加，几何级数给 \(\vartheta(x)\le2Cx+O(1)\)。
\end{proof}
这个界没有给渐近常数一，却足够控制以后 Laplace 变换中的函数。若先用素数定理推出该界，再把它放进素数定理的证明，就会形成循环。

\subsection{\texorpdfstring{\(\psi(x)\)}{\AnalysisPdfPsi(x)}}
\label{b3:subsec:270103}

定义 \(\psi(x)=\sum_{n\le x}\Lambda(n)=\sum_{p^r\le x}\log p\)。按幂次分组，
\[
\psi(x)=\sum_{r\ge1}\vartheta(x^{1/r}),\qquad
0\le\psi(x)-\vartheta(x)\le C\sqrt x\log(x+2).
\]
因为仅有 \(r\le\log_2x\) 的项，每个 \(r\ge2\) 的项由 Chebyshev 界不超过 \(C\sqrt x\)。于是 \(\psi(x)=O(x)\)，并且高次素数幂相对于 \(x\) 是低阶项。
\symidx{\(\pi(x)\)：素数计数函数}
\symidx{\(\vartheta(x)\)：素数对数权和}
\symidx{\(\psi(x)\)：素数幂对数权和}

\subsection{von Mangoldt 函数}
\label{b3:subsec:270104}

前面的上界可以与另一个初等二项式估计结合，得到正确的增长数量级。中心二项式系数是 \((1+1)^{2n}\) 展开中的最大项，所以
\[
 \binom{2n}{n}\ge\frac{4^n}{2n+1}.
\]
它的素数 \(p\) 指数为
\[
 \sum_{j\ge1}\left(\left\lfloor\frac{2n}{p^j}\right\rfloor
                   -2\left\lfloor\frac n{p^j}\right\rfloor\right).
\]
每个括号是零或一，并且 \(p^j>2n\) 时为零。因此取对数得到
\[
 \psi(2n)\ge\log\binom{2n}{n}
       \ge2n\log2-\log(2n+1).
\]
对一般大实数 \(x\)，取 \(n=\lfloor x/2\rfloor\) 并用单调性，便知存在常数 \(c_0>0\)，使 \(\psi(x)\ge c_0x\)。扣去高次幂的 \(o(x)\) 误差，\(\vartheta(x)\) 也具有同样形式的下界。

\ProseParagraphBreak

于是对充分大的 \(x\)，存在正常数 \(c_1,C_1\)，使
\[
 c_1\frac{x}{\log x}\le\pi(x)\le C_1\frac{x}{\log x}.
\]
下界由 \(\vartheta(x)\le\pi(x)\log x\) 得到。对于上界，先分开 \(p\le\sqrt x\) 的至多 \(\sqrt x\) 个素数；其余每个 \(\log p\ge\tfrac12\log x\)，故个数不超过 \(2\vartheta(x)/\log x\)。这说明初等估计已经找到正确尺度。素数定理增加的内容是把上下界的两个不同常数收紧到同一个精确常数一。

\ProseParagraphBreak

von Mangoldt 函数的非负性使 \(\psi\) 单调。这一点将在 Tauber 步骤中把平均信息变成点态渐近；符号改变的系数即使有相同的解析极点，也可能有很大的振荡。这里的 \(\Lambda\) 是\secref{b3:sec:2501}的算术函数，不是上一章的完备 zeta 函数。

\subsection{基本渐近等价式}
\label{b3:subsec:270105}

\begin{proposition}\label{b3:prop:27-equivalences}
三个断言 \(\psi(x)\sim x\)、\(\vartheta(x)\sim x\)、\(\pi(x)\sim x/\log x\) 两两等价。
\end{proposition}
\begin{proof}
前两个由高次幂误差 \(o(x)\) 等价。若 \(\vartheta(t)\sim t\)，部分求和给
\(\pi(x)=\vartheta(x)/\log x+\int_2^x\vartheta(t)/(t\log^2t)\,\mathrm{d}t\)。由 \(\vartheta(t)=O(t)\)，将积分在 \(\sqrt x\) 分开，可得它为 \(O(\sqrt x)+O(x/\log^2x)=o(x/\log x)\)，于是得到 \(\pi\) 的渐近。
反之若 \(\pi(x)\sim x/\log x\)，则 \(\vartheta(x)=\pi(x)\log x-\int_2^x\pi(t)/t\,\mathrm{d}t\)。\(\pi(t)=O(t/\log t)\)，同样在 \(\sqrt x\) 分割得积分 \(O(\sqrt x\log x)+O(x/\log x)=o(x)\)，所以 \(\vartheta(x)\sim x\)。
\end{proof}

% !TeX root = ../../Book3.tex
\section{\texorpdfstring{\(\zeta\)}{\AnalysisPdfZeta} 函数与素数}
\label{b3:sec:2702}

带非负系数的对数导数把 \(\psi\) 转成一个解析函数。极点将提供主项，边界零点则可能产生与主项同量级的振荡，所以二者要分别处理。

\subsection{对数导数}
\label{b3:subsec:270201}

前章与本章采用同一个等式
\[
F(s)\coloneqq-\zeta'(s)/\zeta(s)
=\sum_{n\ge1}\Lambda(n)n^{-s}
=s\int_1^\infty\psi(x)x^{-s-1}\,\mathrm{d}x,
\quad\operatorname{Re}s>1.
\]
最后一步是 Abel 求和，端点 \(\psi(x)x^{-s}\to0\) 由 \(\psi(x)=O(x)\) 保证。由于 \(\zeta(s)=1/(s-1)+\gamma+O(s-1)\)，有 \(F(s)=1/(s-1)-\gamma+O(s-1)\)。

\subsection{Euler 乘积}
\label{b3:subsec:270202}

在 \(\sigma>1\)，Euler 乘积给
\[
\operatorname{Re}F(\sigma+\mathrm{i}t)
=\sum_{p,r\ge1}(\log p)p^{-r\sigma}\cos(rt\log p).
\]
单个余弦项会改变符号；把三个高度的值按合适系数组合后，则可使每项非负。绝对收敛允许先在每个素数幂处组合，再求全和。

\subsection{显式公式观点}
\label{b3:subsec:270203}

从形式上看，\(F(s)x^s/s\) 在 \(s=1\) 的留数为 \(x\)，而在零点 \(\rho\) 处的留数为 \(-m_\rho x^\rho/\rho\)。这说明零点实部为什么影响误差幂次。本章的证明采用 Tauber 方法，无需先建立一般反演显式公式；这段留数计算只解释单个奇点的局部贡献，不将无穷多留数未经截断相加。

% !TeX root = ../../Book3.tex
\section{\texorpdfstring{\(\operatorname{Re}s=1\)}{Re(s)=1} 上的无零点性}
\label{b3:sec:2703}

无零点性只需要一个非负三角多项式和局部 Laurent 展开。这里不会证明一条定量宽度的无零区域；所证内容恰是边界直线上每个有限高度没有零点。

\subsection{Hadamard–de la Vallée Poussin 论证}
\label{b3:subsec:270301}

以下经典论证称为\term[hadamard-de-la-vallee-poussin]{Hadamard--de la Vallée Poussin 无零点论证}[Hadamard--de la Vallée Poussin zero-free argument]。设 \(t\ne0\)，若 \(1+\mathrm{i}t\) 是 \(\zeta\) 的 \(m\) 阶零点，则
\(F(s)=-m/(s-1-\mathrm{i}t)+O(1)\) 于该点附近。若 \(1+2\mathrm{i}t\) 也是零点，记阶数为 \(m_2\)，否则记零。

\subsection{正性技巧}
\label{b3:subsec:270302}

恒等式 \(3+4\cos u+\cos2u=2(1+\cos u)^2\ge0\) 给
\begin{equation}\label{b3:eq:27-positivity}
3F(\sigma)+4\operatorname{Re}F(\sigma+\mathrm{i}t)
+\operatorname{Re}F(\sigma+2\mathrm{i}t)\ge0,\quad\sigma>1.
\end{equation}
每个权 \(\Lambda(n)n^{-\sigma}\) 非负，故相乘求和保持不等式。把中间系数改小会丧失排除单零点所需的优势；三个系数的具体数值承担了证明任务。

\subsection{边界估计}
\label{b3:subsec:270303}

令 \(\sigma=1+\varepsilon\)。局部展开给
\(\varepsilon F(1+\varepsilon)\to1\)，
\(\varepsilon\operatorname{Re}F(1+\varepsilon+\mathrm{i}t)\to-m\)，第三项极限为 \(-m_2\)。固定 \(t\) 使每个展开只在一个固定邻域使用，无需关于所有高度的一致估计。

\subsection{零点排除}
\label{b3:subsec:270304}

\begin{theorem}\label{b3:thm:27-zero-free-line}
Riemann zeta 函数在直线 \(\operatorname{Re}s=1\) 上没有零点；\(s=1\) 处仍为单极点。
\end{theorem}
\begin{proof}
若 \(t\ne0\) 处有 \(m\ge1\) 阶零点，将\eqref{b3:eq:27-positivity}乘 \(\varepsilon\) 并取极限，得到 \(3-4m-m_2\ge0\)，但右式不超过负一，矛盾。\(t=0\) 处已知为极点。
\end{proof}
所以 \(F(s)-1/(s-1)\) 在整条闭半平面 \(\operatorname{Re}s\ge1\) 的一个开邻域中全纯；“一个开邻域”允许向左伸出的宽度随高度改变，不暗含统一条带。下一节使用的正是这种局部延拓。

\ProseParagraphBreak

这里可以逐点说明这个邻域如何获得。在不等于一的边界点，\(\zeta\) 全纯且非零，连续性使它在一个足够小的圆盘内仍非零，因而对数导数全纯。在一处，写
\[
 \zeta(s)=\frac{a(s)}{s-1},\qquad a(1)=1,
\]
其中 \(a\) 在小圆盘内全纯非零，则
\[
 F(s)-\frac1{s-1}=-\frac{a'(s)}{a(s)}
\]
在该圆盘全纯。把这些圆盘与原右半平面取并，便得到需要的开邻域。对于每一段有限高度的边界，紧致性还能选出统一的正宽度；当高度增长时，这个宽度却可能缩小。

\ProseParagraphBreak

另一个等价的正性检验从 \(\log\zeta\) 出发。在 \(\sigma>1\) 内逐项应用同一三角多项式并取指数，得
\[
 \zeta(\sigma)^3|\zeta(\sigma+\mathrm it)|^4
              |\zeta(\sigma+2\mathrm it)|\ge1.
\]
若 \(1+\mathrm it\) 是 \(m\ge1\) 阶零点，左侧在 \(\sigma\downarrow1\) 时为
\(O((\sigma-1)^{4m+m_2-3})\to0\)，矛盾。指数 \(4m+m_2-3\) 与前面留数不等式的符号相反，两者实际计算了同一局部阶数。这个版本也再次显示为何需要系数四压过极点贡献三。

% !TeX root = ../../Book3.tex
\section{Tauber 型方法}
\label{b3:sec:2704}

现在已经知道，\(-\zeta'/\zeta\) 在 \(s=1\) 的主部为 \(1/(s-1)\)，并且扣除这一项后可以全纯越过整条直线 \(\operatorname{Re}s=1\)。尚未得到的是 \(\psi(x)/x\to1\)：变换把不同位置的质量混合在一起，边界信息不能未经解释就变成每一点的计数渐近。本节把这一步单独证明。我们先说明累积量怎样产生变换，再用一个完整的 Tauber 论证，把解析信息还原为原函数的渐近。

\subsection{Abel 型与 Tauber 型定理}
\label{b3:subsec:270401}

素数幂的权重是离散的，但以下方法只使用质量非负和累积量单调，不要求质量集中在整数上。因此先允许一个一般函数 \(A:[1,\infty)\to[0,\infty)\)，假设它右连续且非递减，并约定 \(A(1^-)=0\)。这个约定的意思是：在开始计数以前没有质量，而 \(A(1)\) 可以不为零。若 \(A\) 在一点跳跃，该跳跃就是该点的质量；不能把 \(\mathrm dA\) 一律理解成 \(A'(x)\,\mathrm dx\)。

\begin{proposition}[累积量与 Stieltjes 变换]
\label{b3:prop:27-stieltjes-transform}
设 \(A:[1,\infty)\to[0,\infty)\) 右连续且非递减，每个 \(A(x)\) 都有限，并令 \(A(1^-)=0\)。在 \([1,\infty)\) 上存在唯一的局部有限正 Borel 测度 \(\mu_A\)，满足
\[
 \mu_A([1,x])=A(x)\qquad(x\ge1).
\]
记 \(\mathrm dA=\mathrm d\mu_A\)。对每个实数 \(\sigma>0\)，有允许两边同为无穷大的恒等式
\begin{equation}
\label{b3:eq:27-stieltjes-tonelli}
 \int_{[1,\infty)}x^{-\sigma}\,\mathrm dA(x)
 =\sigma\int_1^\infty A(x)x^{-\sigma-1}\,\mathrm dx.
\end{equation}
若左端对每个 \(\sigma>1\) 有限，则
\[
 F(s)=\int_{1^-}^\infty x^{-s}\,\mathrm dA(x)
 \qquad(\operatorname{Re}s>1)
\]
绝对收敛且全纯，并满足
\begin{equation}
\label{b3:eq:27-stieltjes-transform}
 \frac{F(s)}s=\int_1^\infty A(x)x^{-s-1}\,\mathrm dx,
 \qquad\operatorname{Re}s>1.
\end{equation}
其中 \(1^-\) 表示积分包含 \(1\) 处的原子。
\end{proposition}
\begin{proof}
先构造测度，避免把有跳跃的函数误当成连续函数。令
\(M=\sup_{x\ge1}A(x)\in[0,\infty]\)。若 \(M=0\)，取零测度。若 \(M>0\)，对 \(0<y<M\) 令
\[
 q(y)=\inf\{x\ge1:A(x)\ge y\}.
\]
因为 \(y<M\)，集合非空而且下确界有限。单调性保证它是一个向右延伸的区间；取其中趋于下确界的点列，右连续性给出 \(A(q(y))\ge y\)。因此下确界实际属于这个集合，且
\[
 q(y)\le x\quad\Longleftrightarrow\quad y\le A(x).
\]
这也证明 \(q\) 可测。把 \((0,M)\) 上的 Lebesgue 测度经 \(q\) 推前，所得正 Borel 测度满足
\[
 \mu_A([1,x])
 =\bigl|\{y\in(0,M):y\le A(x)\}\bigr|=A(x).
\]
当 \(M<\infty\) 且 \(A(x)=M\) 时，右端仍为 \(M\)；排除单个端点 \(M\) 不改变 Lebesgue 测度。每个 \([1,x]\) 的质量有限，故 \(\mu_A\) 局部有限。区间质量给出
\[
 \mu_A(\{1\})=A(1),\qquad
 \mu_A((a,b])=A(b)-A(a)\quad(1\le a<b).
\]
这些区间生成 Borel 集，在每个有限区间上用有限测度的唯一性，再取递增并，便得到唯一性。

下面建立积分公式。对 \(y\ge1\) 和实数 \(\sigma>0\)，有
\[
 y^{-\sigma}=\sigma\int_y^\infty x^{-\sigma-1}\,\mathrm dx.
\]
被积函数非负，Tonelli 定理允许交换积分，即使尚不知道有限性：
\[
 \int y^{-\sigma}\,\mathrm d\mu_A(y)
 =\sigma\int_1^\infty
    \mu_A([1,x])x^{-\sigma-1}\,\mathrm dx.
\]
这就是\eqref{b3:eq:27-stieltjes-tonelli}。\(y=x\) 的端点是否包含不影响外层 Lebesgue 积分，但内层采用 \([1,x]\) 正好与右连续的 \(A(x)\) 一致。

若 \(s=\sigma+\mathrm i\tau\)、\(\sigma>1\)，同样有
\(y^{-s}=s\int_y^\infty x^{-s-1}\,\mathrm dx\)。相应二重积分的绝对值由
\[
 |s|\int_1^\infty A(x)x^{-\sigma-1}\,\mathrm dx
 =\frac{|s|}{\sigma}F(\sigma)<\infty
\]
控制，Fubini 定理遂给出\eqref{b3:eq:27-stieltjes-transform}。最后，在 \(\operatorname{Re}s>1\) 的任一紧集上，取 \(\delta>0\) 使 \(\operatorname{Re}s\ge1+\delta\)。对任意固定整数 \(k\ge0\)，
\[
 (\log x)^k x^{-\operatorname{Re}s}
 \le C_{k,\delta}x^{-1-\delta/2}\qquad(x\ge1).
\]
右端对 \(\mu_A\) 可积，故可在积分号下对 \(s\) 求导，证明了全纯性。
\end{proof}

例如，若 \(A(x)=\sum_{n\le x}a_n\)、\(a_n\ge0\)，则 \(\mu_A=\sum_{n\ge1}a_n\delta_n\)，变换就是 \(\sum a_nn^{-s}\)。若 \(A(x)=cx\)、\(c\ge0\)，测度则是 \(c\delta_1+c\,\mathrm dx\)（连续部分限于 \((1,\infty)\)），所以
\[
 F(s)=c+\frac{c}{s-1}=\frac{cs}{s-1}.
\]
这一个例子提醒我们：\(A(x)/x\to c\) 对应极点的系数 \(c\)，不要求整个变换恰好等于 \(c/(s-1)\)；\(1\) 处的质量和有限区间上的改动会贡献全纯项。

\begin{proposition}[Abel 型正向极限]
\label{b3:prop:27-abelian-limit}
设 \(A:[1,\infty)\to[0,\infty)\) 右连续且非递减，令 \(A(1^-)=0\)。若对某个 \(c\ge0\) 有 \(A(x)/x\to c\)，则其 Stieltjes 变换 \(F(s)=\int_{1^-}^\infty x^{-s}\,\mathrm dA(x)\) 在 \(\operatorname{Re}s>1\) 绝对收敛，并且
\[
 \lim_{\sigma\downarrow1}(\sigma-1)F(\sigma)=c,
 \qquad \sigma\in\mathbb R.
\]
\end{proposition}
\begin{proof}
假设和有限区间上的单调性给出 \(A(x)\le Cx\)。由\eqref{b3:eq:27-stieltjes-tonelli}，变换在右半平面绝对收敛。令
\[
 h(t)=e^{-t}A(e^t)\quad(t\ge0),\qquad
 \sigma=1+\varepsilon\quad(\varepsilon>0).
\]
在\eqref{b3:eq:27-stieltjes-transform}中作 \(x=e^t\) 的代换，得到
\[
 \frac{\varepsilon F(1+\varepsilon)}{1+\varepsilon}
 =\varepsilon\int_0^\infty h(t)e^{-\varepsilon t}\,\mathrm dt.
\]
这里 \(h(t)\to c\)，而权函数 \(\varepsilon e^{-\varepsilon t}\) 在正半轴的积分恰为 \(1\)。给定 \(\eta>0\)，选 \(T_0\) 使 \(|h(t)-c|\le\eta\) 于 \(t\ge T_0\)。于是
\[
 \left|\varepsilon\int_0^\infty
       (h(t)-c)e^{-\varepsilon t}\,\mathrm dt\right|
 \le\varepsilon\int_0^{T_0}|h(t)-c|\,\mathrm dt+\eta.
\]
先令 \(\varepsilon\downarrow0\)，再令 \(\eta\downarrow0\)，得到结论。
\end{proof}

这种由原函数渐近推向变换极限的结果称为 Abel 型结果。反过来，由变换的信息恢复原函数渐近，称为\term[tauber-xing-dingli]{Tauber 型定理}[Tauberian theorem]。正向命题只给出实轴上的极限，甚至没有声称 \(F(s)-c/(s-1)\) 全纯越过 \(s=1\)。下面的反向定理将要求更充分的边界信息，并用非负性防止振荡藏在平均之下。

\ProseParagraphBreak

证明在对数坐标中进行，因为乘法尺度 \(x\mapsto xe^v\) 在 \(t=\log x\) 中变成平移。下表整理随后会同时出现的对象；其中 \(h\) 和 \(r\) 在负半轴补为零。
\begin{table}[htbp]
\centering
\caption{Tauber 证明中的变量和变换}
\label{b3:tab:27-tauber-objects}
\begin{tabularx}{\textwidth}{@{}lX@{}}
\toprule
对象 & 定义与作用 \\
\midrule
\(x,t\) & \(x\ge1\)、\(t=\log x\ge0\)：乘法尺度改为加法尺度。 \\
\(A(x)\) & 非负、右连续、非递减的累积量；目标为 \(A(x)/x\to c\)。 \\
\(h(t)\) & \(h(t)=e^{-t}A(e^t)\)：目标改写为 \(h(t)\to c\)。 \\
\(r(t)\) & \(r(t)=h(t)-c\mathbf1_{[0,\infty)}(t)\)：从原函数中减去主项。 \\
\(F(s)\) & \(F(s)=\int_{1^-}^{\infty}x^{-s}\,\mathrm dA(x)\)，最初在 \(\operatorname{Re}s>1\) 定义。 \\
\(G(z)\) & \(G(z)=F(1+z)/(1+z)-c/z\)：\(r\) 的 Laplace 变换，原边界移到虚轴。 \\
\bottomrule
\end{tabularx}
\end{table}

\(A\) 单调而 \(h\) 未必单调，但对 \(t,v\ge0\)，
\begin{equation}
\label{b3:eq:27-slow-decrease}
 h(t+v)=e^{-t-v}A(e^{t+v})\ge e^{-v}h(t).
\end{equation}
所以一个高值不能立即向右大幅下降。反向使用同一不等式，当 \(-\delta\le v\le0\)、\(t\ge\delta\) 时，有 \(h(t+v)\le e^\delta h(t)\)：一个低值也会限制左侧一小段的高度。这两条单边约束将把平均值的极限变成点态极限。

\subsection{Wiener–Ikehara 定理的解析边界版本}
\label{b3:subsec:270402}

\begin{theorem}[Wiener--Ikehara：解析边界版本]
\label{b3:thm:27-wiener-ikehara}
设 \(A:[1,\infty)\to[0,\infty)\) 为有限值、右连续、非递减函数，令 \(A(1^-)=0\)，并以 \(\mu_A([1,x])=A(x)\) 定义其 Stieltjes 测度。假设
\[
 F(s)=\int_{1^-}^\infty x^{-s}\,\mathrm dA(x)
\]
在 \(\operatorname{Re}s>1\) 绝对收敛。若对某个 \(c\ge0\)，函数 \(F(s)-c/(s-1)\) 全纯延拓到闭半平面 \(\operatorname{Re}s\ge1\) 的一个开邻域，则
\[
 \frac{A(x)}x\longrightarrow c\qquad(x\to\infty).
\]
当 \(c>0\) 时，这等价于 \(A(x)\sim cx\)；当 \(c=0\) 时，结论写作 \(A(x)=o(x)\)。
\end{theorem}

这是\term[wiener-ikehara-dingli]{Wiener--Ikehara 定理}[Wiener--Ikehara theorem]的解析边界版本。开邻域可以在不同高度具有不同宽度，不要求一条宽度固定的无零区域。证明也只会一次使用一个有限高度区间上的边界延拓。

\ProseParagraphBreak

直接恢复 \(h(T)\) 需要同时控制无限多频率。为只使用有限段边界，先用一个频率支集紧的平滑核取平均；为保留\eqref{b3:eq:27-slow-decrease}的不等号，核又必须非负。普通的尖锐频率截断不能满足后一要求，其逆变换会改变符号。下面先对光滑截断作逆 Fourier 变换，再取模平方，明确构造出兼具非负性、快速衰减和紧频率支集的核。

\begin{lemma}[非负的紧频率核]
\label{b3:lem:27-positive-bandlimited-kernel}
存在 \(\phi\in\mathcal S(\mathbb R)\) 和 \(q\in C_c^\infty(\mathbb R)\)，使
\[
 \phi(t)=\int_{\mathbb R}q(\xi)e^{2\pi\mathrm it\xi}\,\mathrm d\xi,
 \qquad \phi\ge0,\qquad\int_{\mathbb R}\phi(t)\,\mathrm dt=1,
\]
且 \(\phi\) 在原点附近严格为正。对每个 \(a>0\)、\(b\in\mathbb R\)，平移缩放核
\[
 \phi_{a,b}(v)=a\phi(a(v-b))
\]
以及它的反射 \(v\mapsto\phi_{a,b}(-v)\)，仍具有非负性、单位积分和光滑紧支集的频率表示。
\end{lemma}
\begin{proof}
取非零、非负的 \(\chi\in C_c^\infty(\mathbb R)\)，置
\[
 v(t)=\int_{\mathbb R}\chi(\xi)e^{2\pi\mathrm it\xi}\,\mathrm d\xi,
 \qquad K=\int_{\mathbb R}|v(t)|^2\,\mathrm dt,
 \qquad\phi(t)=|v(t)|^2/K.
\]
对 \(v\) 求任意阶导数只在被积函数中增加 \((2\pi\mathrm i\xi)^k\)；再对 \(\xi\) 反复分部积分，紧支集使端点项为零，证明所有这些导数任意阶快速衰减。因此 \(v\in\mathcal S(\mathbb R)\)，\(K<\infty\)。又 \(v(0)=\int\chi>0\)，所以 \(K>0\)，并且 \(\phi\) 在原点附近严格为正。非负性、Schwartz 性及积分为 \(1\) 都由定义得到。

展开平方，对两个有限频率区间使用 Fubini 定理，并作变量替换，得到
\[
 \phi(t)=\int_{\mathbb R}q(\xi)e^{2\pi\mathrm it\xi}\,\mathrm d\xi,
 \qquad
 q(\xi)=\frac1K\int_{\mathbb R}
       \chi(\xi+\eta)\overline{\chi(\eta)}\,\mathrm d\eta.
\]
若 \(\operatorname{supp}\chi\subset[-R,R]\)，则 \(q\) 支于 \([-2R,2R]\)；积分号下求导给出光滑性。因此这个表示直接由绝对收敛的积分推出，不需要把非可积函数放进 Fourier 反演公式。

换元给出 \(\int\phi_{a,b}=1\)，频率表示中的函数则变为
\[
 q_{a,b}(\xi)=e^{-2\pi\mathrm ib\xi}q(\xi/a).
\]
它仍光滑且紧支集。反射只把 \(q_{a,b}(\xi)\) 换成 \(q_{a,b}(-\xi)\)，其余性质同样保持。
\end{proof}

这里的归一化与全书 Fourier 变换中的 \(e^{-2\pi\mathrm it\xi}\) 一致。因此下一引理中的 Laplace 边界点是 \(2\pi\mathrm i\xi\)，不是 \(\mathrm i\xi\)。这个因子只改变频率的参数，不改变边界是整条虚轴这一事实。

\begin{lemma}[解析边界给出固定核平均]
\label{b3:lem:27-boundary-average}
设可测函数 \(r:\mathbb R\to\mathbb C\) 在负半轴为零，且对某些常数 \(C>0\)、整数 \(m\ge0\) 满足 \(|r(u)|\le C(1+u)^m\) 于 \(u\ge0\)。设其 Laplace 变换
\[
 G(z)=\int_0^\infty r(u)e^{-zu}\,\mathrm du,
 \qquad\operatorname{Re}z>0,
\]
全纯延拓到虚轴的一个开邻域。若 \(q\in C_c^\infty(\mathbb R)\)，且
\(\phi(t)=\int q(\xi)e^{2\pi\mathrm it\xi}\,\mathrm d\xi\)，则
\[
 (r*\phi)(T)=\int_{\mathbb R}r(u)\phi(T-u)\,\mathrm du
 \longrightarrow0\qquad(T\to+\infty).
\]
\end{lemma}
\begin{proof}
先给正半轴上的积分加上阻尼 \(e^{-\varepsilon u}\)，其中 \(\varepsilon>0\)。由于
\[
 \int_0^\infty |r(u)|e^{-\varepsilon u}\,\mathrm du
 \cdot\int_{\mathbb R}|q(\xi)|\,\mathrm d\xi<\infty,
\]
Fubini 定理给出
\begin{equation}
\label{b3:eq:27-damped-average}
 \begin{aligned}
 &\int_0^\infty r(u)e^{-\varepsilon u}\phi(T-u)\,\mathrm du\\
 &\qquad=\int_{\mathbb R}q(\xi)G(\varepsilon+2\pi\mathrm i\xi)
                         e^{2\pi\mathrm iT\xi}\,\mathrm d\xi.
 \end{aligned}
\end{equation}
此时我们先固定 \(T\)，令 \(\varepsilon\downarrow0\)。左端的控制函数可以取 \(C(1+u)^m|\phi(T-u)|\)，它可积，因为 \(\phi\) 是 Schwartz 函数。右端只涉及固定紧集 \(\operatorname{supp}q\)。这个紧集对应的虚轴线段位于延拓域内，紧致性给出一条足够薄的共同邻域；在该邻域中 \(G\) 连续且一致有界。因此右端一致趋于边界积分，得到
\[
 (r*\phi)(T)=\int_{\mathbb R}
       q(\xi)G(2\pi\mathrm i\xi)e^{2\pi\mathrm iT\xi}\,\mathrm d\xi.
\]
振幅 \(Q(\xi)=q(\xi)G(2\pi\mathrm i\xi)\) 光滑且紧支集。对 \(T>0\) 分部积分一次，端点项为零，故
\[
 |(r*\phi)(T)|
 \le\frac1{2\pi T}\int_{\mathbb R}|Q'(\xi)|\,\mathrm d\xi
 \longrightarrow0.
\]
这一步的常数依赖所选核和所用边界线段；它没有给出所有缩放核共同的误差界。
\end{proof}

最后还需要一条实变量引理。它解释为什么先证明 \(h\) 有界，再把核集中到单侧小区间：具有紧频率支集的核在空间中仍带有尾部，不能把尾部直接丢掉。

\begin{lemma}[由单调累积量撤去平均]
\label{b3:lem:27-monotone-unsmoothing}
设 \(h:\mathbb R\to[0,\infty)\) 可测，在负半轴为零，且对某些常数 \(C>0\)、整数 \(m\ge0\) 满足 \(h(t)\le C(1+t)^m\) 于 \(t\ge0\)。假设
\[
 h(t+v)\ge e^{-v}h(t)\qquad(t,v\ge0).
\]
设 \(\phi\in\mathcal S(\mathbb R)\) 非负，\(\int_{\mathbb R}\phi=1\)，且 \(\phi\) 在原点附近严格为正，记 \(\phi_{a,b}(v)=a\phi(a(v-b))\)。若对某个 \(c\ge0\)，每个固定的 \(a>0\)、\(b\in\mathbb R\) 都满足
\begin{equation}
\label{b3:eq:27-all-fixed-averages}
 \int_{\mathbb R}h(T+v)\phi_{a,b}(v)\,\mathrm dv
 \longrightarrow c\qquad(T\to\infty),
\end{equation}
则 \(h(T)\to c\)。
\end{lemma}
\begin{proof}
第一步只选一个固定核来取得全局有界性。因 \(\phi\) 在原点附近为正，可取 \(\delta_0>0\) 使
\(m_0=\int_0^{\delta_0}\phi(v)\,\mathrm dv>0\)。非负性和单边不等式给出
\[
 \int_{\mathbb R}h(T+v)\phi(v)\,\mathrm dv
 \ge\int_0^{\delta_0}e^{-\delta_0}h(T)\phi(v)\,\mathrm dv
 =m_0e^{-\delta_0}h(T).
\]
左端趋于有限数 \(c\)，因此 \(h(T)\) 在充分大的 \(T\) 上有界。多项式增长界还保证它在有限区间有界，所以存在常数 \(B<\infty\)，使 \(0\le h(T)\le B\) 对所有 \(T\in\mathbb R\) 成立。

现在固定 \(\delta>0\) 和 \(0<\eta<1\)。在 \(\phi_{a,b}\) 中取 \(b=\delta/2\)。由换元，
\[
 \int_0^\delta\phi_{a,\delta/2}(v)\,\mathrm dv
 =\int_{-a\delta/2}^{a\delta/2}\phi(w)\,\mathrm dw
 \longrightarrow1\qquad(a\to\infty).
\]
故可以选一个固定的 \(a\)，使该质量至少为 \(1-\eta\)。将核记为 \(\phi_+\)。对所有 \(T\ge0\)，
\[
 \int_{\mathbb R}h(T+v)\phi_+(v)\,\mathrm dv
 \ge(1-\eta)e^{-\delta}h(T).
\]
沿使 \(h(T)\) 趋于上极限的子列使用\eqref{b3:eq:27-all-fixed-averages}，得到
\begin{equation}
\label{b3:eq:27-tauber-upper}
 \limsup_{T\to\infty}h(T)
 \le\frac{ce^\delta}{1-\eta}.
\end{equation}
这里丢弃目标区间以外的积分，是因为 \(h\) 和核都非负。

再取 \(b=-\delta/2\)，把另一个固定核 \(\phi_-\) 的质量至少 \(1-\eta\) 集中在 \((-\delta,0)\)。对 \(T\ge\delta\)、\(-\delta\le v\le0\)，单边条件应用于 \(T+v\) 和 \(-v\)，给出
\[
 h(T)\ge e^v h(T+v),\qquad
 h(T+v)\le e^{-v}h(T)\le e^\delta h(T).
\]
区间外则用刚取得的全局界 \(B\)。因此
\[
 \int_{\mathbb R}h(T+v)\phi_-(v)\,\mathrm dv
 \le e^\delta h(T)+B\eta.
\]
沿实现下极限的子列取极限，得到
\begin{equation}
\label{b3:eq:27-tauber-lower}
 \liminf_{T\to\infty}h(T)\ge e^{-\delta}(c-B\eta).
\end{equation}
\(B\) 在选择 \(\delta,\eta\) 以前已经固定。于是先令 \(\eta\downarrow0\)，再令 \(\delta\downarrow0\)，上下界合为 \(h(T)\to c\)。若 \(c=0\)，\eqref{b3:eq:27-tauber-upper}和 \(h\ge0\) 已直接给出结论。
\end{proof}

\begin{proof}[\cref{b3:thm:27-wiener-ikehara}的证明]
先由正性取得一个粗增长界，这一步不预先假设 \(A(x)=O(x)\)。对 \(T\ge1\)，在 \([1,e^T]\) 上有 \(x^{-1-1/T}\ge e^{-T-1}\)，所以
\[
 F(1+1/T)\ge e^{-T-1}\mu_A([1,e^T])
 =e^{-T-1}A(e^T).
\]
去极点后的函数在 \(s=1\) 附近全纯，故
\(F(1+1/T)=cT+O(1)\)。因此
\[
 h(T)=e^{-T}A(e^T)\le eF(1+1/T)=O(T+1).
\]
在有限区间上，\(A\) 的单调性也给出有界性。将 \(h\) 在负半轴补为零，便得到适用于整条正半轴的线性增长界。

由\cref{b3:prop:27-stieltjes-transform}及 \(x=e^t\) 的代换，对于 \(\operatorname{Re}z>0\)，有
\[
 \int_0^\infty h(t)e^{-zt}\,\mathrm dt
 =\int_1^\infty A(x)x^{-z-2}\,\mathrm dx
 =\frac{F(1+z)}{1+z}.
\]
这里的等式已经由绝对可积的 Fubini 论证建立，不需要对有跳跃的 \(A\) 求导。设
\[
 r(t)=h(t)-c\mathbf1_{[0,\infty)}(t).
\]
它的 Laplace 变换为
\begin{equation}
\label{b3:eq:27-laplace-subtraction}
 \begin{aligned}
 G(z)&=\frac{F(1+z)}{1+z}-\frac cz\\
 &=\frac{F(1+z)-c/z-c}{1+z}.
 \end{aligned}
\end{equation}
最后一个表达式说明零点处的主部确实抵消：\(F(1+z)-c/z\) 由假设全纯越过虚轴，而分母 \(1+z\) 在虚轴上处处非零。因此 \(G\) 全纯越过整条虚轴。式中另一个可能的分母零点 \(z=-1\) 不在虚轴上，可以从延拓域中去掉，不影响所需邻域。

任取\cref{b3:lem:27-positive-bandlimited-kernel}给出的 \(\phi\)，并固定 \(a>0\)、\(b\in\mathbb R\)。将\cref{b3:lem:27-boundary-average}用于反射核 \(v\mapsto\phi_{a,b}(-v)\)，可得
\[
 \int_{\mathbb R}r(T+v)\phi_{a,b}(v)\,\mathrm dv\longrightarrow0.
\]
注意反射正好把卷积中的 \(T-v\) 换成当前需要的 \(T+v\)。常数部分满足
\[
 \int_{\mathbb R}\mathbf1_{[0,\infty)}(T+v)
       \phi_{a,b}(v)\,\mathrm dv
 =\int_{-T}^\infty\phi_{a,b}(v)\,\mathrm dv\longrightarrow1.
\]
相加即得所有固定核的平均极限\eqref{b3:eq:27-all-fixed-averages}。\(A\) 的单调性给出\eqref{b3:eq:27-slow-decrease}，所以\cref{b3:lem:27-monotone-unsmoothing}推出 \(h(T)\to c\)。最后换回 \(x=e^T\)，得到 \(A(x)/x\to c\)。
\end{proof}

证明中的三个极限有明确次序：对每个固定核和固定 \(T\)，先撤去阻尼 \(\varepsilon\downarrow0\)；核固定以后再令 \(T\to\infty\)；取得全局界以后，才改变核，使其集中在越来越短的单侧区间。核越集中，频率支集越大，而假设没有提供所有高度共同的边界估计。因此这条证明自然给出定性极限，没有自动给出有效误差率。

\ProseParagraphBreak

正性也承担了三个不同任务：它先把变换的实轴大小变成 \(h(T)=O(T+1)\)，再把一个固定的核平均变成全局有界性，最后允许丢弃或控制区间外的质量。如果跳过有界性而直接缩窄核，尾部估计只能使用随 \(T\) 增长的粗界，固定的微小尾质量仍可能产生不可忽略的误差。更弱边界假设下的 Tauber 定理需要额外理论，本节只证明所陈述的解析边界版本。

\begin{remark}[仅有一个边界点的信息不足]
\label{b3:rem:27-one-boundary-point}
取 \(0<\varepsilon<1/\sqrt2\)，令
\[
 A(x)=x\bigl(1+\varepsilon\sin(\log x)\bigr),\qquad x\ge1,
 \qquad A(1^-)=0.
\]
因为 \(1-\varepsilon>0\)，它非负；又
\[
 A'(x)=1+\varepsilon\sin\log x+\varepsilon\cos\log x
 \ge1-\sqrt2\varepsilon>0,
\]
所以它严格递增，但 \(A(x)/x\) 始终振荡。其变换由\eqref{b3:eq:27-stieltjes-transform}给出：
\[
 \begin{aligned}
 F(s)
 &=s\int_1^\infty x^{-s}\bigl(1+\varepsilon\sin\log x\bigr)\,\mathrm dx\\
 &=\frac{s}{s-1}+\frac{\varepsilon s}{(s-1)^2+1},
 \qquad\operatorname{Re}s>1.
 \end{aligned}
\]
第二项的计算只需在 \(x=e^t\) 后使用
\(\sin t=(e^{\mathrm it}-e^{-\mathrm it})/(2\mathrm i)\)，分别积分两个指数函数。因而
\[
 F(s)-\frac1{s-1}
 =1+\frac{\varepsilon s}{(s-1)^2+1}
\]
在 \(s=1\) 附近全纯，却在 \(1\pm\mathrm i\) 有极点。这个例子保留了非负性、单调性和线性增长，只缺少整条边界的延拓，仍不能得到渐近。这里 \(A(1)=1\) 产生的原子已经包含在变换中；若遗漏它，计算还会错误地少掉一个常数项。
\end{remark}

\subsection{对 \texorpdfstring{\(\psi(x)\)}{\AnalysisPdfPsi(x)} 的应用}
\label{b3:subsec:270403}

现在逐项核对定理对素数幂计数的适用性。取
\[
 A(x)=\psi(x)=\sum_{n\le x}\Lambda(n),
 \qquad \mu_A=\sum_{n\ge1}\Lambda(n)\delta_n.
\]
由于 \(\Lambda(n)\ge0\)，\(A\) 非负且非递减。有限区间中只有有限个整数，所以 \(A\) 取有限值；在整数处采用 \(n\le x\) 的约定，使阶梯函数右连续。又 \(\Lambda(1)=0\)，这里 \(A(1)=0\)，因而 \(1\) 处恰好没有原子。

\ProseParagraphBreak

在 \(\operatorname{Re}s>1\) 内，Euler 乘积的对数导数已经给出绝对收敛表示
\[
 F(s)=\int_{1^-}^\infty x^{-s}\,\mathrm dA(x)
 =\sum_{n\ge1}\frac{\Lambda(n)}{n^s}
 =-\frac{\zeta'(s)}{\zeta(s)}.
\]
\cref{b3:thm:27-zero-free-line}排除了边界上除 \(s=1\) 以外的零点；\(s=1\) 处的简单极点则使 \(F(s)-1/(s-1)\) 全纯。\secref{b3:sec:2703}已把这些局部邻域与原右半平面接起来，得到包含整个闭半平面的开延拓域。因此\cref{b3:thm:27-wiener-ikehara}中可以取 \(c=1\)，得到
\begin{equation}
\label{b3:eq:27-psi-asymptotic}
 \psi(x)=x+o(x).
\end{equation}

这一应用没有把 \(\psi(x)=O(x)\) 当作 Tauber 定理的前提；一般定理在自己的证明中取得所需有界性。我们此前证明的 Chebyshev 界仍有独立用途，末节的 Newman 路线会直接使用它。对当前路线而言，绝对收敛、边界解析性和系数非负已经足够。下一节再通过三种计数函数之间的比较，把\eqref{b3:eq:27-psi-asymptotic}换成通常的素数定理。

% !TeX root = ../../Book3.tex
\section{素数定理}
\label{b3:sec:2705}

Tauber 步骤给出的是素数幂的对数权和。现在把这个结果代入本章开头已经证明的等价关系，就能得到通常叙述的素数定理。本节还说明它允许怎样改变计数方式，以及哪些看似相近的定量结论仍未由此证明。

\subsection{素数定理及其加权推论}
\label{b3:subsec:270501}

\begin{theorem}[素数定理]\label{b3:thm:27-pnt}
对实数 \(x\ge2\)，令 \(\pi(x)=\#\{\,p\le x\mid p\text{ 为素数}\,\}\)、\(\vartheta(x)=\sum_{p\le x}\log p\)、\(\psi(x)=\sum_{p^m\le x,\ m\ge1}\log p\)，其中 \(p\) 遍历素数、\(m\) 为正整数。则当 \(x\to\infty\) 时，
\[
 \psi(x)\sim x,
 \quad \vartheta(x)\sim x,
 \quad \pi(x)\sim\frac{x}{\log x}.
\]
\end{theorem}
\begin{proof}
在\subsecref{b3:subsec:270403}中，已经核对了 \(A(x)=\psi(x)\) 右连续、非负、非递减，其 Stieltjes 变换在 \(\operatorname{Re}s>1\) 等于 \(-\zeta'(s)/\zeta(s)\)，且减去 \(1/(s-1)\) 后全纯越过整条边界直线。\cref{b3:thm:27-wiener-ikehara}取 \(c=1\)，给出 \(\psi(x)/x\to1\)。再应用\cref{b3:prop:27-equivalences}，就得到其余两个渐近式；其中去掉高次素数幂只损失 \(O(\sqrt x\log x)=o(x)\)，去掉对数权的积分余项也已在该命题中估计。
\end{proof}

这个结果称为\term[sushu-dingli]{素数定理}[prime number theorem]。它断言相对误差趋零，却没有给出趋零的数值速率。即使知道 \(\pi(x)\) 的整体主项，也不能把两个很接近的端点相减后，仍无条件保留同样的相对精度。

\begin{corollary}[固定比例区间内的素数]
\label{b3:cor:27-fixed-ratio-interval}
设 \(\lambda>1\) 为固定实数，\(\pi(x)\) 为不超过实数 \(x\ge1\) 的素数个数。则当 \(x\to\infty\) 时，
\[
 \#\{\,p\mid x<p\le\lambda x,\ p\text{ 为素数}\,\}
 =\pi(\lambda x)-\pi(x)
 \sim(\lambda-1)\frac{x}{\log x}.
\]
特别地，对这个固定的 \(\lambda\)，存在 \(X_\lambda\)，使每个实数 \(x\ge X_\lambda\) 的区间 \((x,\lambda x]\) 都含有素数。
\end{corollary}
\begin{proof}
素数定理给出
\[
 \begin{aligned}
 \pi(\lambda x)-\pi(x)
 &=\frac{\lambda x}{\log(\lambda x)}-\frac{x}{\log x}
      +o\!\left(\frac{x}{\log x}\right)\\
 &=\frac{x}{\log x}
    \left(\frac{\lambda\log x}{\log x+\log\lambda}-1+o(1)\right).
 \end{aligned}
\]
因为 \(\lambda\) 固定，括号趋于正数 \(\lambda-1\)，证明了渐近式。右侧趋于无穷，所以区间内的素数个数最终为正。
\end{proof}

若让 \(\lambda=\lambda(x)\downarrow1\)，预期主项会缩小到 \((\lambda(x)-1)x/\log x\)，而我们目前只有两个端点误差各为 \(o(x/\log x)\)。这不足以保证误差比新的主项小。因此这里的“固定比例”是推论实际使用的条件；任意短区间的计数还需要额外估计。

\ProseParagraphBreak

去掉或加入一个平滑权重仍由部分求和完成。下面把主项中的常数也算出，以免把“乘上 \(p^\alpha\)”误解为简单地把 \(\pi(x)\) 乘上 \(x^\alpha\)。

\begin{proposition}[素数的幂权和]
\label{b3:prop:27-prime-power-weights}
设 \(\alpha>-1\) 为固定实数。对于实数 \(x\ge2\)，对不超过 \(x\) 的素数求和，则
\[
 \sum_{p\le x}p^\alpha
 \sim\frac{x^{\alpha+1}}{(\alpha+1)\log x}
 \quad(x\to\infty).
\]
渐近关系中的参数 \(\alpha\) 保持固定，不主张对趋近于 \(-1\) 的参数一致成立。
\end{proposition}
\begin{proof}
先计算所需的积分尺度。置 \(\beta=\alpha+1>0\)，并令
\[
 I_\alpha(x)=\int_2^x\frac{t^\alpha}{\log t}\,\mathrm dt,
 \quad D_\alpha(x)=\frac{x^\beta}{\beta\log x}.
\]
这里 \(D_\alpha(x)\to\infty\)，且
\[
 D_\alpha'(x)=\frac{x^\alpha}{\log x}
       \left(1-\frac1{\beta\log x}\right).
\]
给定 \(0<\varepsilon<1\)，当 \(x\ge X\) 且 \(X\) 足够大时，
\((1-\varepsilon)I_\alpha'(x)\le D_\alpha'(x)\le I_\alpha'(x)\)。从 \(X\) 到 \(x\) 积分，固定下端的常数除以 \(D_\alpha(x)\) 后消失，于是
\[
 1\le\liminf_{x\to\infty}\frac{I_\alpha(x)}{D_\alpha(x)}
 \le\limsup_{x\to\infty}\frac{I_\alpha(x)}{D_\alpha(x)}
 \le\frac1{1-\varepsilon}.
\]
令 \(\varepsilon\downarrow0\)，得到 \(I_\alpha(x)\sim D_\alpha(x)\)。这个论证同时适用于正的和负的 \(\alpha\)，只使用了 \(\beta>0\)。

记 \(r(t)=\pi(t)-t/\log t=o(t/\log t)\)。由\cref{b3:prop:25-integral-summation}，
\[
 \sum_{p\le x}p^\alpha
 =x^\alpha\pi(x)-\alpha\int_2^x t^{\alpha-1}\pi(t)\,\mathrm dt.
\]
边界项中的误差 \(x^\alpha r(x)\) 为 \(o(x^\beta/\log x)\)。对积分误差，选 \(X\) 使 \(|r(t)|\le\varepsilon t/\log t\) 对所有 \(t\ge X\) 成立，则
\[
 \left|\int_2^x t^{\alpha-1}r(t)\,\mathrm dt\right|
 \le C_{\alpha,X}+\varepsilon I_\alpha(x).
\]
除以 \(D_\alpha(x)\)，先令 \(x\to\infty\)，再令 \(\varepsilon\downarrow0\)，可知该误差为 \(o(D_\alpha(x))\)。这里初始区间只给固定常数；\(D_\alpha(x)\to\infty\) 正是它能够被忽略的原因。因此
\[
 \begin{aligned}
 \sum_{p\le x}p^\alpha
 &=\frac{x^\beta}{\log x}-\alpha I_\alpha(x)
           +o\!\left(\frac{x^\beta}{\log x}\right)\\
 &=\left(1-\frac\alpha{\alpha+1}\right)
      \frac{x^\beta}{\log x}
           +o\!\left(\frac{x^\beta}{\log x}\right),
 \end{aligned}
\]
即为所求。
\end{proof}

例如，取 \(\alpha=1\) 得 \(\sum_{p\le x}p\sim x^2/(2\log x)\)；取 \(\alpha=-1/2\) 得 \(\sum_{p\le x}p^{-1/2}\sim2\sqrt x/\log x\)。第二个例子也说明，权重虽在衰减，累积和仍可增长。临界值 \(\alpha=-1\) 则必须单独处理，不能把它代入分母 \(\alpha+1\)。

\begin{proposition}[素数倒数和的主项]
\label{b3:prop:27-prime-reciprocal-asymptotic}
对于实数 \(x\ge2\)，对不超过 \(x\) 的素数求和，则当 \(x\to\infty\) 时，
\[
 \sum_{p\le x}\frac1p=\log\log x+o(\log\log x).
\]
\end{proposition}
\begin{proof}
部分求和给出
\[
 \sum_{p\le x}\frac1p
 =\frac{\pi(x)}x+\int_2^x\frac{\pi(t)}{t^2}\,\mathrm dt.
\]
边界项为 \(O(1/\log x)\)。写 \(\pi(t)=t/\log t+r(t)\)，其中 \(r(t)=o(t/\log t)\)，则积分主项为
\[
 \int_2^x\frac{\mathrm dt}{t\log t}
 =\log\log x-\log\log2.
\]
给定 \(\varepsilon>0\)，取固定 \(X>2\)，使 \(|r(t)|\le\varepsilon t/\log t\) 对所有 \(t\ge X\) 成立。于是
\[
 \left|\int_2^x\frac{r(t)}{t^2}\,\mathrm dt\right|
 \le C_X+\varepsilon(\log\log x-\log\log X).
\]
除以 \(\log\log x\)，先令 \(x\to\infty\)，再令 \(\varepsilon\downarrow0\)，得到所需余项。
\end{proof}

这个结果给出了素数倒数和发散的主项。误差 \(o(\log\log x)\) 本身仍可能无界；要判断 \(\sum_{p\le x}1/p-\log\log x\) 是否收敛到常数，需要进一步控制误差积分。

\ProseParagraphBreak

素数定理还有两个常用的表述：给定素数的序号，估计这个素数有多大；或者用对数积分作为比较函数。两者都是主项层面的改写，下面分别说明如何恢复其中的对数和误差尺度。

\begin{corollary}[第 \texorpdfstring{\(n\)}{n} 个素数的渐近]
\label{b3:cor:27-nth-prime}
设 \(p_n\) 为按递增顺序排列的第 \(n\) 个素数，\(n\ge1\)。则
\[
 p_n\sim n\log n\quad(n\to\infty).
\]
\end{corollary}
\begin{proof}
素数定理保证 \(\pi(x)\to\infty\)，所以每个 \(p_n\) 均存在，且 \(p_n\to\infty\)。在 \(x=p_n\) 处有精确等式 \(\pi(p_n)=n\)，因而
\[
 n=\frac{p_n}{\log p_n}(1+o(1)).
\]
这里不能先把 \(\log p_n\) 擅自换成 \(\log n\)，因为这正是尚须证明的一步。对上式取实对数，由 \(\log(1+o(1))=o(1)\)，有
\[
 \log n=\log p_n-\log\log p_n+o(1).
\]
除以趋于无穷的 \(\log p_n\)，并使用 \(\log\log p_n/\log p_n\to0\)，得到 \(\log n/\log p_n\to1\)。最后把两个已知比值相乘：
\[
 \frac{p_n}{n\log n}
 =\frac{p_n}{n\log p_n}\cdot\frac{\log p_n}{\log n}
 \longrightarrow1.
\]
\end{proof}

这个推导展示了反求渐近量时的一般步骤：先从计数关系取得含未知量的等式，再证明缓慢变化的对数可以替换，最后代回。它没有把“函数近似于某函数”直接当作“反函数近似于随意代换后的表达式”。

\begin{proposition}[对数积分与主项尺度]
\label{b3:prop:27-logarithmic-integral}
对实数 \(x\ge2\)，定义
\[
 \operatorname{Li}(x)\coloneqq\int_2^x\frac{\mathrm dt}{\log t}.
\]
则当 \(x\to\infty\) 时，
\[
 \operatorname{Li}(x)=\frac{x}{\log x}
          +O\!\left(\frac{x}{\log^2x}\right),
 \quad
 \pi(x)\sim\operatorname{Li}(x),
\]
其中 \(\pi(x)\) 是不超过 \(x\) 的素数个数。
\end{proposition}
\begin{proof}
由 \((t/\log t)'=1/\log t-1/\log^2t\)，分部积分给
\[
 \operatorname{Li}(x)
 =\frac{x}{\log x}-\frac2{\log2}
             +\int_2^x\frac{\mathrm dt}{\log^2t}.
\]
当 \(x\ge4\) 时，把余项分在 \(\sqrt x\)，有
\[
 0\le\int_2^x\frac{\mathrm dt}{\log^2t}
 \le\frac{\sqrt x}{\log^22}+\frac{4x}{\log^2x}
 =O\!\left(\frac{x}{\log^2x}\right).
\]
最后一步使用 \(\log^2x/\sqrt x\to0\)。于是 \(\operatorname{Li}(x)\sim x/\log x\)；再与素数定理比较，得到 \(\pi(x)/\operatorname{Li}(x)\to1\)。这个比值只在足够大的 \(x\) 上使用，因而 \(\operatorname{Li}(2)=0\) 不造成问题。
\end{proof}

这里的\term[duishu-jifen]{对数积分}[logarithmic integral]从二开始积分，不涉及一处的奇点；采用其他常用归一化时可能相差一个常数，本章固定使用上面的定义。由当前证明，我们只可断言
\[
 \pi(x)-\operatorname{Li}(x)=o\!\left(\frac{x}{\log x}\right).
\]
\(\operatorname{Li}(x)\) 与 \(x/\log x\) 之间的差有更小的上界，并不意味着 \(\pi(x)-\operatorname{Li}(x)\) 自动也有这个上界。要证明对数积分提供更精细的误差近似，必须另行估计素数计数的余项。

\subsection{误差项与零点的关系}
\label{b3:subsec:270502}

素数定理给出 \(\psi(x)=x+o(x)\)，它刻画了误差相对于主项的大小。若误差能够改进为一个严格小于一的固定幂次，变换就会向左多收敛一段；由此可以严格推出一个无零点半平面。误差积分的绝对收敛将使这一联系具体化。

\begin{proposition}[幂次误差蕴含无零点半平面]
\label{b3:prop:27-power-error-zero-free}
设 \(\zeta\) 为 Riemann zeta 函数的亚纯延拓，且对于实数 \(x\ge1\) 定义 \(\psi(x)=\sum_{p^r\le x,\ r\ge1}\log p\)，其中 \(p\) 遍历素数、\(r\) 为正整数。若存在固定实数 \(0\le\alpha<1\)、\(B\ge0\)，使
\[
 \psi(x)=x+O\bigl(x^\alpha(\log x)^B\bigr)
 \quad(x\to\infty),
\]
则 \(\zeta(s)\) 在严格半平面 \(\operatorname{Re}s>\alpha\) 内没有零点；\(s=1\) 仍为单极点。
\end{proposition}
\begin{proof}
令 \(E(x)=\psi(x)-x\)。由假设并增大常数，可对所有 \(x\ge2\) 使用 \(|E(x)|\le Cx^\alpha(\log x)^B\)。在 \(\operatorname{Re}s>\alpha\) 定义
\[
 H(s)=s\int_1^\infty E(x)x^{-s-1}\,\mathrm dx.
\]
积分在 \([1,2]\) 上没有困难，因为 \(E\) 在该紧区间有界。若 \(K\) 是这个半平面内的紧集，令 \(\delta=\min_{s\in K}(\operatorname{Re}s-\alpha)>0\)。则尾部被常数倍的
\[
 x^{-1-\delta}(\log x)^B
\]
一致控制。对参数求任意固定次数 \(j\) 的导数，只会再出现至多 \((\log x)^j\) 的因子；代换 \(u=\log x\) 后，控制函数的积分化为常数倍的
\[
 \int_{\log2}^\infty e^{-\delta u}u^{B+j}\,\mathrm du<\infty.
\]
故参数积分在紧集上局部一致收敛，且可逐项求导，\(H\) 在 \(\operatorname{Re}s>\alpha\) 全纯。这也说明固定对数幂不会妨碍严格半平面内的延拓。

在原来的 \(\operatorname{Re}s>1\) 中，由\secref{b3:sec:2702}的部分求和表示，
\[
 -\frac{\zeta'(s)}{\zeta(s)}
 =s\int_1^\infty\psi(x)x^{-s-1}\,\mathrm dx
 =\frac{s}{s-1}+H(s).
\]
注意主项为 \(s/(s-1)\)，因为 \(s\int_1^\infty x^{-s}\,\mathrm dx=s/(s-1)\)。它与 \(1/(s-1)\) 只差全纯常数一，但在精确的积分恒等式中应保留这个因子。

左右两侧都是连通半平面 \(\operatorname{Re}s>\alpha\) 上的亚纯函数，并在其非空开子集 \(\operatorname{Re}s>1\) 相等。亚纯函数的恒等定理于是使等式在整个半平面成立。右侧除 \(s=1\) 外全纯。

若 \(\rho\ne1\) 是该半平面内 \(\zeta\) 的 \(m\ge1\) 阶零点，局部可写成
\[
 \zeta(s)=(s-\rho)^m u(s),\quad u(\rho)\ne0,
\]
其中 \(u\) 全纯，并在足够小的圆盘内非零。于是
\[
 -\frac{\zeta'(s)}{\zeta(s)}
 =-\frac m{s-\rho}-\frac{u'(s)}{u(s)}.
\]
其极点留数为非零整数 \(-m\)，不可能可去，和右侧在 \(\rho\) 全纯矛盾。\(s=1\) 已知是极点，不是零点，命题得证。
\end{proof}

这里得到的是 \(\operatorname{Re}s>\alpha\)，没有包括边界 \(\operatorname{Re}s=\alpha\)。在边界上，刚才的控制积分失去了因子 \(e^{-\delta u}\)，原误差假设不再保证收敛。特别地，即使误差是 \(O(x^{1/2}(\log x)^B)\)，这个方向也只直接排除实部严格大于二分之一的零点。

\ProseParagraphBreak

反向从无零点区域推出定量误差，还需要控制 \(\zeta'/\zeta\) 的大小、反演的截断高度和围道边上的积分。仅仅说明某个区域没有零点，还没有完成这些估计。本章的 Tauber 证明按“先固定边界紧段，再取变量趋于无穷”的次序获得 \(o(x)\)；它没有暗含有效的无零区域宽度，也没有自动给出通常更强的误差界。

% !TeX root = ../../Book3.tex
\OptionalSection{Newman 的简化解析证明}{b3:sec:2706}

Newman 的方法用一个带端点消去因子的 Cauchy 积分直接证明 Laplace 变换的 Tauber 结论。它可用素数的一次幂权和完成另一个证明，避免一般的正核论证。

\subsection{解析准备}
\label{b3:subsec:270601}

\begin{theorem}[Newman 解析引理]\label{b3:thm:27-newman}
设 \(f:[0,\infty)\to\mathbb C\) 有界且局部可积，\(g(z)=\int_0^\infty f(t)e^{-zt}\,\mathrm{d}t\) 于 \(\operatorname{Re}z>0\) 定义。若 \(g\) 全纯延拓越过整个虚轴，则 \(\int_0^T f(t)\,\mathrm{d}t\to g(0)\)。
\end{theorem}
这一结果称为\term[newman-jiexi-yinli]{Newman 解析引理}[Newman's analytic lemma]。这里结论为广义积分收敛，不要求绝对可积。以下围道证明可核对 Zagier 的《Newman's Short Proof of the Prime Number Theorem》\cite[pp. 707--708]{Zagier1997Newman}。

\subsection{围道论证}
\label{b3:subsec:270602}

\begin{proof}[\cref{b3:thm:27-newman}的证明]
记 \(|f|\le B\)、\(g_T(z)=\int_0^T f(t)e^{-zt}\,\mathrm{d}t\)。固定 \(R>0\)，由虚轴紧段的延拓可取 \(\delta>0\)，使 \(g\) 在圆盘 \(|z|\le R\) 与 \(\operatorname{Re}z\ge-\delta\) 的交集及边界邻域全纯。沿这一区域正向边界 \(C\)，Cauchy 公式给
\[
g(0)-g_T(0)=\frac1{2\pi\mathrm{i}}\int_C(g(z)-g_T(z))e^{Tz}
\left(1+\frac{z^2}{R^2}\right)\frac{\mathrm{d}z}{z}.
\]
在右半圆 \(z=Re^{\mathrm{i}\theta}\) 上，
\(|g-g_T|\le Be^{-T\operatorname{Re}z}/\operatorname{Re}z\)，而 \(|1+z^2/R^2|=2|\cos\theta|\)，故整个被积函数模不超过 \(2B/R^2\)，积分贡献至多 \(B/R\)。端点处按连续极限理解这一相消后的界。
在左侧路径，将 \(g\) 与 \(g_T\) 分开。因为 \(g_T\) 整且左侧变形不跨零，含 \(g_T\) 的路径可改成完整左半圆。对 \(\operatorname{Re}z<0\)，有
\(|g_T(z)e^{Tz}|\le B/|\operatorname{Re}z|\)，相同因子相消再次给至多 \(B/R\)。剩下的 \(g\) 积分沿固定左路径，其上除两个端点外实部严格负，故 \(e^{Tz}\to0\)。路径不经过零，\(g(z)(1+z^2/R^2)/z\) 连续有界，且 \(|e^{Tz}|\le1\)，支配收敛使这一积分趋零。由此
\(\limsup_{T\to\infty}|g(0)-g_T(0)|\le2B/R\)。任意增大 \(R\) 即得结论；延拓宽度 \(\delta\) 可随 \(R\) 缩小，不影响先取 \(T\) 极限的次序。
\end{proof}

\subsection{初等 Tauber 型归约}
\label{b3:subsec:270603}

若非递减函数 \(A(x)\) 满足 \(\int_1^\infty(A(x)-x)x^{-2}\,\mathrm{d}x\) 收敛，则 \(A(x)/x\to1\)。若存在任意大 \(x\) 使 \(A(x)\ge\lambda x\)、\(\lambda>1\)，则在 \([x,\lambda x]\) 上
\[
\int_x^{\lambda x}\frac{A(t)-t}{t^2}\,\mathrm{d}t
\ge\int_1^\lambda\frac{\lambda-u}{u^2}\,\mathrm{d}u>0,
\]
与收敛积分尾段趋零矛盾。若 \(A(x)\le\lambda x\)、\(0<\lambda<1\)，则在 \([\lambda x,x]\) 上用 \(A(t)\le A(x)\)，得到一个固定严格负上界，同样矛盾。这一步只使用单调性，并不要求 \(A\) 连续或可微。

\ProseParagraphBreak

两个固定差距可以分别算成
\[
 \int_1^\lambda\frac{\lambda-u}{u^2}\,\mathrm du
       =\lambda-1-\log\lambda>0\quad(\lambda>1),
\]
以及
\[
 \int_\lambda^1\frac{\lambda-u}{u^2}\,\mathrm du
       =1-\lambda+\log\lambda<0\quad(0<\lambda<1).
\]
严格符号来自 \(\log\lambda<\lambda-1\)。积分收敛意味着任意两个同时趋于无穷的端点之间，带符号积分也趋零；所以只要比值 \(A(x)/x\) 与一保持一个固定距离，单调性就会把这一次偏离延续成一个乘法区间上的固定积分偏差。这是这里由平均恢复逐点渐近的全部实变量内容。

\subsection{素数定理的另一种证明}
\label{b3:subsec:270604}

令 \(P(s)=\sum_p(\log p)p^{-s}\)。高次幂部分
\(H(s)=\sum_{p,r\ge2}(\log p)p^{-rs}\) 在 \(\operatorname{Re}s>1/2\) 局部一致绝对收敛，故 \(P=F-H\) 在直线一上除一处主部外全纯延拓。Abel 求和给
\[
\frac{P(1+z)}{1+z}-\frac1z
=\int_0^\infty\bigl(e^{-t}\vartheta(e^t)-1\bigr)e^{-zt}\,\mathrm{d}t.
\]
括号内函数由 Chebyshev 界有界；左侧在虚轴全纯，因为一处的单极点已抵消，其他点由无零点性保证。Newman 引理遂给 \(\int_1^\infty(\vartheta(x)-x)x^{-2}\,\mathrm{d}x\) 收敛。上一小节推出 \(\vartheta(x)\sim x\)，再去掉对数权即得素数定理。此证法同样不能仅凭无零点性给出定量误差，其围道半径是先固定、最后才趋无穷的。


\begin{chaptersummary}
素数定理的解析证明先选择合适的权重。Euler 对数导数对应非负的 von Mangoldt 函数；高次素数幂的贡献相对 \(x\) 为低阶项，部分求和再去掉对数权。Chebyshev 估计独立给出正确数量级，不能从最终定理借来充当前提。

边界无零点性来自三个高度之间的正性约束。局部 Laurent 展开把一个假设零点变成负留数，系数 \(3,4,1\) 使它无法与正性相容。所得结论保证去主部后的变换在每个有限高度范围有一个左侧邻域，却不提供所有高度共用的宽度。

Wiener--Ikehara 证明先用紧频率正核得到平均极限，再由非负性取得有界性，最后以单调夹逼恢复逐点渐近。Newman 证明则用消去因子控制半圆端点，只把截断整函数的路径移入较大的左半圆；变换本身始终留在已知延拓域内。二者都对每个固定尺度先取极限，所以只得到定性主项。固定比例区间、第 \(n\) 个素数和若干加权渐近可以由此推出；变化的短区间、精确次项及有效误差仍需新的估计。
\end{chaptersummary}

\BookExercises

% 八题自拟或按现行正文重新组织，不虚构外部题号。
% 原第n素数、Abel正向结论、单点边界反例及Li主项已在正文证明。
% 题7为选读；其余题只使用本章主线与已建立的积分/级数工具。

\begin{exercise}\label{b3:ex:27-1}
设 \(\lambda>1\) 为固定实数，求和遍历素数。证明
\[
 \sum_{x<p\le\lambda x}\frac1p
       \sim\frac{\log\lambda}{\log x}\qquad(x\to\infty).
\]
说明为何不能直接相减两个仅有 \(o(\log\log x)\) 误差的素数倒数和渐近式。
\end{exercise}
\begin{solution}
由\eqref{b3:eq:25-interval-summation}，含上端、不含下端的区间和为
\[
 \frac{\pi(\lambda x)}{\lambda x}-\frac{\pi(x)}x
       +\int_x^{\lambda x}\frac{\pi(t)}{t^2}\,\mathrm dt.
\]
写 \(\pi(t)=t(1+\varepsilon(t))/\log t\)，其中 \(\varepsilon(t)\to0\)。令
\(e_x=\sup_{t\ge x}|\varepsilon(t)|\)，则 \(e_x\to0\)。端点项的差等于
\[
 \frac1{\log(\lambda x)}-\frac1{\log x}
       +O\!\left(\frac{e_x}{\log x}\right)
 =O((\log x)^{-2})+o((\log x)^{-1}).
\]
积分中的误差绝对值不超过
\(e_x\int_x^{\lambda x}(t\log t)^{-1}\,\mathrm dt\)，因此积分为
\[
 (1+o(1))\log\frac{\log(\lambda x)}{\log x}
       \sim\frac{\log\lambda}{\log x}.
\]
最后一步使用固定的 \(\log\lambda>0\) 和 \(\log(1+u)\sim u\)。端点误差相对这个主项趋零，结论成立。

若只相减 \(\sum_{p\le y}p^{-1}=\log\log y+o(\log\log y)\)，原有误差尺度远大于 \(1/\log x\)，不能保证差值的相对误差消失。这里重新在目标区间上部分求和，才取得与问题相配的尺度。
\end{solution}

\begin{exercise}\label{b3:ex:27-2}
令 \(p_n\) 为第 \(n\) 个素数。证明
\[
 \frac{p_{n+1}}{p_n}\longrightarrow1,\qquad
 p_{n+1}-p_n=o(p_n),\qquad
 \sum_{k=1}^n p_k\sim\frac12n^2\log n.
\]
解释这些结论为何没有给出每个相邻素数间隔为 \(O(\log p_n)\)。
\end{exercise}
\begin{solution}
由\cref{b3:cor:27-nth-prime}，\(p_n=n\log n(1+o(1))\)。于是
\[
 \frac{p_{n+1}}{p_n}
 =\frac{(n+1)\log(n+1)}{n\log n}\cdot
       \frac{1+o(1)}{1+o(1)}\longrightarrow1.
\]
减一再乘 \(p_n\)，即得间隔为 \(o(p_n)\)。

为处理前 \(n\) 个素数的和，先由单调函数 \(t\log t\) 的上下积分比较得到
\[
 \sum_{k=2}^n k\log k
 =\frac12n^2\log n+O(n^2).
\]
这里原函数为 \(\tfrac12t^2\log t-\tfrac14t^2\)；和与积分之差至多为末项的一个常数倍，\(O(n\log n)\) 已可被上面的误差吸收。给定 \(\epsilon>0\)，取 \(K\) 使 \(k\ge K\) 时
\[
 (1-\epsilon)k\log k\le p_k\le(1+\epsilon)k\log k.
\]
从 \(K\) 加到 \(n\)，有限头部除以 \(n^2\log n\) 趋零。先令 \(n\to\infty\)，再令 \(\epsilon\downarrow0\)，便得到系数 \(1/2\)。

\(o(p_n)\) 允许的误差比 \(\log p_n\) 大得多。把两个仅知相对误差趋零的 \(p_n\) 渐近式相减，也不能把误差自动缩小到对数尺度。因此这里没有得到对每一个间隔的那种上界。
\end{solution}

\begin{exercise}\label{b3:ex:27-3}
对实数 \(x\ge2\)，定义有限乘积
\[
 M(x)=\prod_{p\le x}(1-p^{-1})^{-1}.
\]
证明
\[
 \log M(x)-\sum_{p\le x}\frac1p
 \longrightarrow C_0\coloneqq\sum_p\sum_{r\ge2}\frac1{rp^r}<\infty,
 \qquad M(x)=(\log x)^{1+o(1)}.
\]
仅由本题所得误差，能否进一步断言 \(M(x)/\log x\) 收敛到某个正常数？
\end{exercise}
\begin{solution}
每个因子为正，有限乘积可以取实对数。由 \(-\log(1-u)=\sum_{r\ge1}u^r/r\)、\(0<u<1\)，有准确等式
\[
 \log M(x)=\sum_{p\le x}\frac1p
          +\sum_{p\le x}\sum_{r\ge2}\frac1{rp^r}.
\]
尾部非负，且
\[
 \sum_p\sum_{r\ge2}\frac1{rp^r}
 \le\sum_p\frac1{p(p-1)}
 \le\sum_{n=2}^\infty\frac1{n(n-1)}=1.
\]
故随 \(x\) 增大的第二项单调趋于有限值 \(C_0\)。由\cref{b3:prop:27-prime-reciprocal-asymptotic}，
\[
 \log M(x)=\log\log x+o(\log\log x).
\]
除以 \(\log\log x\) 后取指数，即为
\(M(x)=\exp((1+o(1))\log\log x)=(\log x)^{1+o(1)}\)。

要推出 \(M(x)/\log x\) 有正常数极限，需要控制
\(\log M(x)-\log\log x\) 到收敛常数的尺度。当前只知道它为 \(o(\log\log x)\)，仍可能无界或振荡。\(C_0\) 的存在只解决高次幂部分，尚未控制 \(\sum_{p\le x}p^{-1}-\log\log x\)。这不是否认更强结果，而是区分本题实际证明的内容。
\end{solution}

\begin{exercise}\label{b3:ex:27-4}
令 \(\Lambda\) 为 von Mangoldt 函数，并定义非负算术函数
\[
 a(n)=\sum_{d\mid n}\frac{\Lambda(d)}d,\qquad
 A(x)=\sum_{n\le x}a(n)\quad(x\ge1).
\]
用本章的 Wiener--Ikehara 定理证明
\[
 A(x)\sim c x,\qquad
 c=\sum_{d\ge1}\frac{\Lambda(d)}{d^2}
   =-\frac{\zeta'(2)}{\zeta(2)}>0.
\]
再用有限求和与支配收敛独立核验这个渐近。
\end{exercise}
\begin{solution}
在 \(\operatorname{Re}s>1\)，绝对收敛允许把约数和写成 \(n=dk\)，得到
\[
 D_a(s)=\sum_{n\ge1}\frac{a(n)}{n^s}
 =\zeta(s)\sum_{d\ge1}\frac{\Lambda(d)}{d^{s+1}}
 =\zeta(s)F(s+1),
 \qquad F=-\frac{\zeta'}{\zeta}.
\]
绝对收敛的依据是
\(\zeta(\sigma)\sum_d\Lambda(d)d^{-\sigma-1}<\infty\)，\(\sigma>1\)。函数 \(F(s+1)\) 在 \(\operatorname{Re}s>0\) 全纯，因为 \(s+1\) 位于 Euler 乘积已经保证无零点的半平面。故 \(D_a\) 在实部正的半平面只有一处可能的极点：一处的单极点，留数为 \(F(2)=c\)。因此
\(D_a(s)-c/(s-1)\) 在这一半平面全纯。又 \(c>0\)，例如 \(d=2\) 的一项已严格为正。

\(A\) 非负、右连续且非递减，其 Stieltjes 测度是
\(\sum_n a(n)\delta_n\)，变换恰为 \(D_a\)。\cref{b3:thm:27-wiener-ikehara}的全部条件满足，故 \(A(x)/x\to c\)。

另一证明完全不需要边界理论。交换有限和可得
\[
 \frac{A(x)}x
 =\sum_{d\le x}\frac{\Lambda(d)}d\,\frac{\lfloor x/d\rfloor}{x}.
\]
将 \(d>x\) 的项补为零。对每个固定 \(d\)，该项趋于 \(\Lambda(d)/d^2\)，且总由同一个可和数列 \(\Lambda(d)/d^2\) 控制。因此可交换极限与求和，得到同一常数。这说明 Tauber 定理能统一处理此题，但此例本身已有更直接的绝对可和尾界，并不依赖素数定理。
\end{solution}

\begin{exercise}\label{b3:ex:27-5}
下面构造说明：若丢掉单调性，即使变换的去主部部分能延拓到整个复平面，也不能只凭非负性和线性增长推出原函数渐近。令
\[
 \phi(u)=\max(1-|u|,0),\qquad
 \delta_n=\tfrac14e^{-n^2}\quad(n\ge1),
\]
并对 \(t\ge0\)、\(x\ge1\) 定义
\[
 h(t)=1+\sum_{n\ge1}\phi\!\left(\frac{t-n}{\delta_n}\right),
 \qquad A(x)=x h(\log x).
\]
证明 \(A\) 连续、非负且为 \(O(x)\)，但非单调，且 \(A(x)/x\) 无极限。再证明
\[
 F(s)=s\int_1^\infty A(x)x^{-s-1}\,\mathrm dx,
 \qquad \operatorname{Re}s>1,
\]
满足 \(F(s)-1/(s-1)\) 有整函数延拓。
\end{exercise}
\begin{solution}
第 \(n\) 个峰仅支撑在 \([n-\delta_n,n+\delta_n]\)。这些区间互不相交，而且每个紧区间只与有限多个相交。因此 \(h\) 连续且 \(1\le h\le2\)，\(A\) 连续、非负且 \(x\le A(x)\le2x\)。在整数和半整数处，
\[
 h(n)=2,\qquad h(n+\tfrac12)=1,
\]
所以 \(A(x)/x=h(\log x)\) 无极限。又
\[
 A(e^n)=2e^n,\qquad A(e^{n+\delta_n})=e^{n+\delta_n}<2e^n,
\]
因为 \(\delta_n\le1/4<\log2\)。后一点较大而函数值较小，故 \(A\) 不是非递减函数。

当 \(\operatorname{Re}s>1\) 时，由 \(A(x)\le2x\)，原积分绝对收敛。令 \(x=e^t\)，有
\[
 F(s)=\frac{s}{s-1}+s\sum_{n\ge1}B_n(s),
 \quad
 B_n(s)=\delta_n e^{-(s-1)n}
     \int_{-1}^1\phi(u)e^{-(s-1)\delta_nu}\,\mathrm du.
\]
每个 \(B_n\) 是整函数。若紧集上 \(|s-1|\le L\)，则
\[
 |B_n(s)|\le2\delta_n e^{L(n+\delta_n)}.
\]
右边可和，因为 \(\delta_n=\tfrac14e^{-n^2}\) 的下降压过任意固定指数 \(e^{Ln}\)。所以 \(\sum B_n\) 局部一致收敛，定义整函数 \(B\)。从而
\[
 F(s)-\frac1{s-1}=1+sB(s)
\]
延拓为整函数。

窄峰的高度始终为一，而总宽度下降得非常快，积分变换因而看见很小的面积。单调累积量却不能这样在一个峰后立即下降；本章正是利用这一限制从平均恢复点态。这里的 \(F\) 是题设给出的积分变换，不能未经单调性就把 \(dA\) 当作非负测度来应用 Tauber 定理。
\end{solution}

\begin{exercise}\label{b3:ex:27-6}
令 \(\vartheta(x)=\sum_{p\le x}\log p\)、
\(\psi(x)=\sum_{p^r\le x,r\ge1}\log p\)。证明更精确的差值渐近
\[
 \psi(x)-\vartheta(x)\sim\sqrt x\qquad(x\to\infty).
\]
在证明中分别指出哪一步只用 Chebyshev 上界，哪一步确实使用素数定理。
\end{exercise}
\begin{solution}
按幂次分组，把平方项单独取出：
\[
 \psi(x)-\vartheta(x)
 =\vartheta(\sqrt x)
       +\sum_{3\le r\le\log_2x}\vartheta(x^{1/r}).
\]
Chebyshev 上界给 \(\vartheta(y)\le Cy\)；对 \(y<2\) 可取其值为零。因此余和非负且至多为
\[
 Cx^{1/3}\log_2x=o(\sqrt x).
\]
这一误差估计没有用到素数定理。素数定理只在识别平方项的精确系数时使用：
\(\vartheta(\sqrt x)\sim\sqrt x\)。将两部分合起来，即得结论。

因此高次素数幂虽然不改变 \(\psi(x)\sim x\) 的首项，却在 \(\psi-\vartheta\) 的较小尺度上留下首项 \(\sqrt x\)。证明主项相同与证明二者之差的渐近，需要不同精度的信息。
\end{solution}

\begin{optionalexercise}\label{b3:ex:27-7}
Newman 解析引理中不能仅把“有界”放宽成 \(O(t)\)。令
\[
 f(t)=t e^{\mathrm it^2}\quad(t\ge0).
\]
证明其 Laplace 变换 \(g(z)=\int_0^\infty f(t)e^{-zt}\,\mathrm dt\) 在 \(\operatorname{Re}z>0\) 绝对收敛，并具有整函数延拓，但 \(\int_0^T f(t)\,\mathrm dt\) 不收敛。可将积分射线旋转到辐角 \(\pi/4\)，但须证明圆弧项消失。
\end{optionalexercise}
\begin{solution}
若 \(\sigma=\operatorname{Re}z>0\)，原积分的绝对值由
\(\int_0^\infty t e^{-\sigma t}\,\mathrm dt<\infty\) 控制。对固定这样的 \(z\)，考虑整函数
\(w e^{\mathrm iw^2-zw}\)，沿扇形 \(0\le\arg w\le\pi/4\)、\(|w|\le R\) 使用 Cauchy 定理。圆弧上 \(w=Re^{\mathrm i\theta}\)，被积函数连同弧长元的模为
\[
 R^2\exp\!\bigl(-R^2\sin2\theta
               -R\operatorname{Re}(ze^{\mathrm i\theta})\bigr)\,\mathrm d\theta.
\]
由于 \(\operatorname{Re}(ze^{\mathrm i\theta})\) 在零处为 \(\sigma>0\)，可以固定一个 \(0<\theta_0<\pi/4\)，使 \(0\le\theta\le\theta_0\) 时它至少为 \(\sigma/2\)。这一段弧积分的模不超过
\((\pi/4)R^2e^{-R\sigma/2}\)，趋零。余下角度上
\(\sin2\theta\ge\sin2\theta_0>0\)，故其模不超过
\[
 \frac\pi4R^2e^{-R^2\sin2\theta_0+R|z|},
\]
也趋零。因此两条射线上的无穷积分相等。令 \(w=e^{\mathrm i\pi/4}u\)，注意 \(w\,\mathrm dw=\mathrm i u\,\mathrm du\)，得到
\[
 g(z)=\mathrm i\int_0^\infty
             u e^{-u^2-z e^{\mathrm i\pi/4}u}\,\mathrm du.
\]
右边对所有 \(z\in\mathbb C\) 都有意义。在任意 \(|z|\le M\) 上，核及其各阶参数导数被常数倍
\(u^{k+1}e^{-u^2+Mu}\) 控制；这在正半轴可积，故右边是整函数，给出所需延拓。特别地 \(g(0)=\mathrm i/2\)。

然而直接求原函数得到
\[
 \int_0^T t e^{\mathrm it^2}\,\mathrm dt
       =\frac{e^{\mathrm iT^2}-1}{2\mathrm i}.
\]
取 \(T=\sqrt{2\pi n}\) 时值为零，取 \(T=\sqrt{(2n+1)\pi}\) 时值为 \(\mathrm i\)，所以不存在极限。此例保留了整个边界的解析延拓，撤去的确实只有有界增长假设。原围道证明中依赖固定上界 \(B\) 的两项 \(B/R\) 估计因而不能照搬。
\end{solution}

\begin{exercise}\label{b3:ex:27-8}
设 \(a>0\) 为固定实数，\(\Lambda\) 为 von Mangoldt 函数。用素数定理证明广义指数平滑的渐近
\[
 S_a(x)\coloneqq\sum_{n\ge1}\Lambda(n)e^{-(n/x)^a}
       \sim\Gamma(1+1/a)\,x\qquad(x\to\infty).
\]
特别地，\(a=1\) 时得到 \(\sum_n\Lambda(n)e^{-n/x}\sim x\)。说明这一步属于由原计数推向变换的 Abel 方向，而非素数定理的另一个独立证明。
\end{exercise}
\begin{solution}
对固定 \(x>0\)，\(\Lambda(n)\le\log n\)，指数 \(e^{-(n/x)^a}\) 压过任意负幂，所以级数绝对收敛。将 \(\psi(t)=\sum_{n\le t}\Lambda(n)\) 在 \(0\le t<1\) 延为零。由正函数积分恒等式
\[
 e^{-(n/x)^a}
   =\int_{n/x}^\infty a u^{a-1}e^{-u^a}\,\mathrm du
\]
和 Tonelli 定理，可把求和移入积分，得到
\[
 \frac{S_a(x)}x
   =\int_0^\infty\frac{\psi(xu)}x\,
                       a u^{a-1}e^{-u^a}\,\mathrm du.
\]
对每个 \(u>0\)，素数定理给 \(\psi(xu)/x\to u\)。由 Chebyshev 上界，存在常数 \(C\)，使所有 \(t\ge0\) 均有 \(\psi(t)\le Ct\)，所以被积函数不超过
\(Ca u^a e^{-u^a}\)。它在原点和无穷端都可积。支配收敛于是给
\[
 \lim_{x\to\infty}\frac{S_a(x)}x
 =\int_0^\infty a u^a e^{-u^a}\,\mathrm du
 =\int_0^\infty v^{1/a}e^{-v}\,\mathrm dv
 =\Gamma(1+1/a),
\]
其中换元为 \(v=u^a\)。Gamma 值正且有限，故极限等价于题设的渐近。\(a=1\) 时常数是 \(\Gamma(2)=1\)。

本证明在逐点极限处已经使用 \(\psi(t)\sim t\)，所以只是把已知计数渐近平滑化。若要从平滑结果倒推原计数，则还需相应的 Tauber 论证，不能倒读上面的支配收敛步骤。
\end{solution}
